% Kapitel 2 -- Von Kuonrâts Geburt und seiner frühen Jugend
\section{Von Kuonrâts Geburt und seiner frühen Jugend}

\mylettrine{U}{m} das Jahr des Herrn zwölfhundert und dreizehn, so wird mir von glaubwürdigen Männern versichert, erblickte ein Knabe das Licht der Welt, dem man den Namen \gls{Kuonrat von Steinare} gab. Er war der erstgeborene Sohn \gls{Hadmar von Steinare}s, geboren in den Mauern eben jener Burg Ottenstein, über welche im vorangehenden Abschnitt berichtet ward; und seine Mutter war eine Frau, deren Name mir nicht überliefert ist, da die Urkunden darüber schweigen und die Zeugen, die ich befragte, sich nicht mehr erinnern konnten oder wollten.

Was von Kuonrâts frühester Kindheit berichtet wird, ist wenig und unzuverlässig; ich will daher nicht erfinden, was ich nicht weiß. Eines jedoch steht fest und ist durch glaubwürdige Berichte belegt: im elften Jahre seines Lebens ward jene Entscheidung getroffen, die seinen künftigen Lebensweg zu formen versprach. Der Herr \gls{Wolfhart von Raabs} zu \gls{Raabs an der Thaya} hatte nur eine einzige Tochter, \gls{Diemut von Raabs} mit Namen. Da er ohne Sohn war und keine Aussicht auf einen männlichen Erben hatte, bestand die Gefahr, dass sein \gls{lehen} nach seinem Tode an den Landesfürsten heimfiele oder in fremde Hände geriet. Der Landesfürst entschied, dass dem nicht so sein solle, und bestimmte, dass der junge \gls{Kuonrat von Steinare} mit der Tochter des Herrn von Raabs verlobt werde, auf dass er nach seiner Mündigkeit das \gls{lehen} zu \gls{Raabs an der Thaya} übernehme und das Geschlecht von Raabs durch seine Heirat fortführe.

So ward \gls{Kuonrat von Steinare}, noch ein Kind von elf Jahren, mit \gls{Diemut von Raabs} verlobt. Da er nun als Erbe von Raabs vorgesehen war, brauchte das Haus von Steinære einen anderen Mann für die väterlichen Güter. Dieser war \gls{Nikolaus von Steinare}.  Er wurde wahrhaftig schlicht Niclâs getauft, als vermöchten ihre Zungen die rechte lateinische Form \enquote{Nicolaus}, die doch einem Heiligen eignet, nicht gar wohl hervorzubringen. Er war der jüngere Bruder Kuonrâts, über den mir wenig berichtet ward, außer dass er dem Vater an Wesen und Statur ähnlicher gewesen sein soll als der ältere Bruder.

Über \gls{Kuonrat von Steinare}s Kindheit in jenen Jahren, von seiner Geburt bis zu seinem vierzehnten Lebensjahre, ist mir nichts Weiteres bekannt. Ob er die Tage auf Burg Ottenstein verbrachte, ob er zu Raabs gebracht wurde, ob er reiten lernte oder lesen, ob er weinte oder lachte, der Herr allein weiß es. Was aus jener Stille ward, davon handeln die folgenden Abschnitte.